%%%%%%%%%%%%%%%%%%%%%%%%%%%%%%%%%%%%%%%%%%%%%%%%%%%%%%%%%%%%%%%%%%%%%%%%%%%%%%%%
%   Copyright 2009, Oracle and/or its affiliates.
%   All rights reserved.
%
%
%   Use is subject to license terms.
%
%   This distribution may include materials developed by third parties.
%
%%%%%%%%%%%%%%%%%%%%%%%%%%%%%%%%%%%%%%%%%%%%%%%%%%%%%%%%%%%%%%%%%%%%%%%%%%%%%%%%

\chapter{Changes Since Fortress 1.0 Specifications}
\applabel{spec-changes}

\note{Changes Since
http://projectfortress.sun.com/Projects/Community/changeset/1444}

\section{Changes Made by the 2026 Revival to the Working Draft}
\seclabel{revival-changes}
\sloppy

This section lists the changes made to this working draft by the 2026
revival of the Fortress project (\url{https://github.com/palimondo/fortress}).
The revival works to finish what the designers intended, judged by this
working draft, and changes the text only where a decision it has taken for
the implementation would otherwise leave the text and the implementation in
open disagreement.
Each changed passage carries a box headed
``Revised by the 2026 revival'' that points to its entry here.

Each entry gives the affected sections, the change, its rationale, its
effect, the original text, and what reversing to the alternative not taken
would require.
The original text is that of the \emph{Working Draft of February 2011},
the unrevised copy of this document kept in the revival's repository under
the directory \nolinkurl{Specification-1.0-frozen/}, cited by path and line
there.
The directory's name is misleading: its sources are the working draft of
2~February 2011, and only the file \nolinkurl{fortress.1.0.pdf} in it is the
Fortress~1.0 release.
The full reasoning behind every entry, with each original passage and each
new one, is in the revival's decision record,
\nolinkurl{explorations/compile-ladder/rung-spec-route-a/decision-record.md}
in the same repository.

Every change below follows from one decision of the revival, taken on
24~September 2026 and called \emph{route~A}: the compiled type checker keeps
the type group's rule that no type is a subtype of two different
instantiations of one parameterized trait, a rule this working draft never
stated, and the library's numeric types are to be reorganized so that they
satisfy it.
The alternative not taken, \emph{route~C}, is described once, in
\secref{revival-routec}.

\subsection{Instantiation exclusion}
\seclabel{revival-rule}

\begin{description}
\item[Affected sections] \secref{trait-types}; \secref{trait-decls}.
\item[Change] A property of the exclusion relation is added:
two instantiations of the same parameterized trait exclude each other
unless every static argument other than an operator argument is the same.
The rule it imposes on declarations is added beside the example of
\TYP{InclusiveMolecule}: a trait or object may not extend two such
instantiations, directly or through its supertraits.
\item[Rationale] Without the rule a value may belong to two instantiations
of one parameterized trait, and code specialised for one of them can be
chosen for a caller that was promised the other; with the rule relaxed, the
compiled implementation accepts such a program and fails at run time.
The type group adopted the rule as \emph{multiple instantiation exclusion}
(Allen, Hilburn, Kilpatrick, Luchangco, Ryu, Chase, and Steele, ``Type
Checking Modular Multiple Dispatch with Parametric Polymorphism and
Multiple Inheritance'', OOPSLA 2011), and the compiled type checker has
enforced it since 2010.
The team's later, unfinished Types chapter (Luchangco, 2012) states it in
the two parts used here, under the name \emph{instantiation exclusion}.
Operator arguments do not count: an operator argument only names the
operator a subtrait inherits (\secref{opr-ident}), neither implementation
distinguishes instantiations by it at run time, and the algebra of the
library chapters depends on types that are several instantiations of one
trait differing only in operator arguments.
\KWD{nat}, \KWD{int} and \KWD{bool} arguments count: they are values at
run time (\secref{natparams}, \secref{boolparams}).
\item[Effect] Programs in which a type is a subtype of two instantiations
of one parameterized trait that differ in any other static argument are no
longer valid. The compiled type checker refuses them where the arguments
that differ are types, or numerals given as \KWD{nat} or \KWD{int}
arguments (row~402 of the revival's gap ledger). It does not yet refuse
them where those arguments are \KWD{bool} arguments, though it accepts boolean parameters written out (row~406),
or where one of them is a static parameter of the trait or object that extends both (row~414).
The interpreter does not check the rule.
The library's numeric chapters are not yet revised
(\secref{revival-notrevised}).
\item[Original text] The Working Draft of February 2011 states no such
property; its relations are ``the smallest ones that satisfy all the
properties given in those sections (and this one)''
(\nolinkurl{Specification-1.0-frozen/basic/types-vals-vars.tex}, lines
185--189), and the properties it gives for trait types are these (lines
212--216):
\begin{quote}
A trait declaration may also include an \KWD{excludes} clause,
in which case the defined trait type
excludes every type that appears in that clause.
Every trait type also excludes every arrow type and every tuple type,
and the special types \TYP{()} and \TYP{BottomType}.
\end{quote}
Its rule on declarations is only that of exclusive traits,
``no trait can extend them both''
(\nolinkurl{Specification-1.0-frozen/basic/traits.tex}, lines 218--222).
\item[Route C] Replaces this rule by the rule of \secref{revival-routec}.
\end{description}

\subsection{Where-clause variables in extends clauses}
\seclabel{revival-where}

\begin{description}
\item[Affected sections] \secref{where-clauses}; \secref{trait-decls}.
\item[Change] One sentence is added: a \KWD{where}-clause variable may not
appear as a static argument in an \KWD{extends} clause.
The three declarations that followed it are kept, marked
``Not allowed'', each with the way to write it
(\secref{revival-covariance}, \secref{revival-every},
\secref{revival-empty}).
A box after Victor Luchangco's two draft notes in \secref{trait-decls},
which say that a trait can explicitly extend itself through hidden type
variables and that ``we could just rule this a special case'', answers
them.
\item[Rationale] Such a variable makes the declared type a subtype of one
instantiation for each type the variable may stand for, and instantiation
exclusion forbids any two of them.
\item[Effect] None of the three declarations ran before the revision:
both implementations stop at name resolution, because a
\KWD{where}-clause variable is not in scope in an \KWD{extends} clause.
The revision makes that the rule.
The calculus of \appref{where-calculus} still admits such
declarations (\secref{revival-notrevised}).
\item[Original text] No such sentence; the draft notes are unchanged.
\item[Route C] The same under route~C, which places no type below
infinitely many instantiations of one trait.
\end{description}

\subsection{The covariance example}
\seclabel{revival-covariance}

\begin{description}
\item[Affected section] \secref{where-clauses}.
\item[Change] The declaration is kept, marked ``Not allowed''; the claims
that trait declarations may extend other instantiations of themselves and
that the declaration expresses covariance are withdrawn; static
parameters are stated to be invariant, and a widening function with
bounded static parameters is given as the way to convert one
instantiation to another.
\item[Rationale] \EXP{C\llbracket{}S\rrbracket} would be a subtype of
\EXP{C\llbracket{}T\rrbracket} for every supertype \VAR{T} of \VAR{S}.
The widening function is how the library writes covariance, in the
conversions and the \OPR{APPCOV} operator of its covariant collections,
and it runs on both implementations.
\item[Effect] Covariance is written as an explicit conversion, with the
static arguments given at the call. The box at the example adds that the
type group's later design declares variance with a \KWD{covariant}
modifier, implemented on neither implementation; the decision record and
row~404 of the revival's gap ledger say what is known of it.
\item[Original text] \nolinkurl{Specification-1.0-frozen/basic/trait-parameters.tex},
lines 339--351; the declaration is the one kept above, marked
``Not allowed'', and the text around it read:
\begin{quote}
Trait declarations are allowed to extend other instantiations of
themselves. For example, we can write:
[the declaration]
In this declaration, for every subtype \VAR{S} of \VAR{T}, \EXP{C\llbracket{}S\rrbracket}
is a subtype of \EXP{C\llbracket{}T\rrbracket}. Effectively, we have expressed
  the fact that the static parameter \VAR{S} of \VAR{C} is covariant.
\end{quote}
\item[Route C] The same under route~C.
\end{description}

\subsection{A subtrait of every instantiation}
\seclabel{revival-every}

\begin{description}
\item[Affected section] \secref{where-clauses}.
\item[Change] ``the following trait declaration is legal'' becomes
``is not allowed''; the paragraph after the declaration is kept, in the
conditional, and a closing sentence says why the declaration is not
allowed.
\item[Rationale] The paragraph is the team's own account of why such a
trait is dangerous, written before the rule; it becomes the reason for the
rule.
\item[Effect] None beyond \secref{revival-where}.
\item[Original text] \nolinkurl{Specification-1.0-frozen/basic/trait-parameters.tex},
lines 354--380; the declaration is the one kept above, and the text around
it read:
\begin{quote}
Trait declarations need not have any static parameters in order
to have a \KWD{where} clause. For example, the following trait
declaration is legal:
[the declaration]
In this declaration, trait \VAR{C} is a subtrait of \emph{every}
instantiation of parametric trait \VAR{D}. Thus, trait \VAR{C} has all
of the methods of every instantiation of \VAR{D}. By thinking of the
declaration this way, we can see what restrictions we need to impose on
the trait \VAR{C} in order for it to be sensible. [The rest of the
paragraph is unchanged.]
\end{quote}
\item[Route C] The same under route~C.
\end{description}

\subsection{The empty list}
\seclabel{revival-empty}

\begin{description}
\item[Affected section] \secref{where-clauses}.
\item[Change] The declaration of \TYP{Empty} is kept, marked
``Not allowed''; \EXP{\TYP{Empty}\llbracket{}T\rrbracket}, one empty list
for each element type, is given beside it.
\item[Rationale] The original makes \TYP{Empty} a
\EXP{\TYP{List}\llbracket{}T\rrbracket} for every \VAR{T} at once.
The replacement is the library's own shape for its empty collections and
the same choice as for \TYP{Nothing} (\secref{revival-nothing}); a program
of this shape runs on both implementations.
\item[Effect] The element type is written at each use, as in
\EXP{\TYP{Empty}\llbracket\mathbb{Z}32\rrbracket}, until the static
arguments of a generic object can be inferred, which neither
implementation does today. The section no longer has a running example of
a \KWD{where} clause on an object.
\item[Original text] \nolinkurl{Specification-1.0-frozen/basic/trait-parameters.tex},
lines 383--399; the declaration is the one kept above, introduced by:
\begin{quote}
Object or functional declarations may include \KWD{where} clauses.
Here is an example declaration of an \TYP{Empty} list:
\end{quote}
\item[Route C] The same under route~C.
\end{description}

\subsection{The empty optional value}
\seclabel{revival-nothing}

\begin{description}
\item[Affected sections] \secref{convenience-types};
\secref{chained-exceptions}; \chapref{lib:exception}.
\item[Change] \TYP{Nothing} becomes
\EXP{\KWD{value}\;\;\KWD{object}\:\TYP{Nothing}\llbracket{}T\rrbracket}
extending \EXP{\TYP{Maybe}\llbracket{}T\rrbracket}, with no
\KWD{excludes} clause, and \TYP{Maybe} comprises
\EXP{\TYP{Nothing}\llbracket{}T\rrbracket} and
\EXP{\TYP{Just}\llbracket{}T\rrbracket}; the original declaration is kept,
marked ``Not allowed''. The defaults of an exception's message and chain
become \EXP{\TYP{Nothing}\llbracket\TYP{String}\rrbracket} and
\EXP{\TYP{Nothing}\llbracket\TYP{Exception}\rrbracket}.
\item[Rationale] The original makes \TYP{Nothing} a
\EXP{\TYP{Maybe}\llbracket{}T\rrbracket} for every \VAR{T} at once. Its
\KWD{excludes} clause is also one the grammar gives no object declaration
(\secref{object-decls}), and is redundant, an object trait type already
excluding every type that is not its supertype (\secref{object-traits}).
The replacement is the library's declaration.
\item[Effect] The empty case is written with its type, as in
\EXP{\TYP{Nothing}\llbracket\TYP{String}\rrbracket}, until the static
arguments of a generic object can be inferred from the type expected.
The compiled path's own library, used until the one library replaces it,
still declares a single \TYP{Nothing} with a coercion to every
\EXP{\TYP{Maybe}\llbracket{}T\rrbracket}, and refuses
\EXP{\TYP{Nothing}\llbracket\TYP{String}\rrbracket}
(row~331 of the revival's gap ledger).
\item[Original text] \nolinkurl{Specification-1.0-frozen/basic-lib/convenience.tex},
lines 34--53: the declaration kept above, marked ``Not allowed'', together
with:
\begin{quote}
An optional value \VAR{v} is either \TYP{Nothing} or
\EXP{\TYP{Just}(v)} declared as follows:
[\ldots]
\begin{Fortress}
\(\KWD{trait} \TYP{Maybe}\llbracket{}T\rrbracket \KWD{comprises} \{\,\TYP{Nothing}, \TYP{Just}\llbracket{}T\rrbracket\,\}\)
\end{Fortress}
[\ldots]
\end{quote}
\nolinkurl{Specification-1.0-frozen/basic/exceptions.tex}, lines 111--130:
\begin{quote}
These fields are default to \TYP{Nothing} as follows: [\ldots]
\begin{Fortress}
{\tt~~}\pushtabs\=\+\(  \KWD{getter} \VAR{message}()\COLON \TYP{Maybe}\llbracket\TYP{String}\rrbracket = \TYP{Nothing}\)\\
\(  \ldots\)\\
\(  \KWD{getter} \VAR{chain}()\COLON \TYP{Maybe}\llbracket\TYP{Exception}\rrbracket = \TYP{Nothing}\)\-\\\poptabs
\end{Fortress}
[\ldots]
where an optional value \VAR{v} is either \TYP{Nothing} or
\EXP{\TYP{Just}(v)}.
\end{quote}
\nolinkurl{Specification-1.0-frozen/basic-lib/exception.tex}, lines 22--24:
\begin{quote}
These fields are default to \TYP{Nothing}
where an optional value \VAR{v} is either \TYP{Nothing} or
\EXP{\TYP{Just}(v)} as declared in \secref{convenience-types}.
\end{quote}
\item[Route C] The same under route~C.
\end{description}

\subsection{The root's algebra}
\seclabel{revival-object}

\begin{description}
\item[Affected sections] \chapref{lib:object}.
\item[Change] \TYP{Object} extends \TYP{Any} alone and \TYP{Tuple} likewise,
each still excluding the other; the sentences saying that they describe
their operators with algebraic constraints are removed. The original
declaration of \TYP{Object} is kept, marked ``Not allowed''.
\item[Rationale] Every trait but \TYP{Any} and \TYP{Object} extends
\TYP{Object}, and every commutative algebraic type \VAR{T} is an
\EXP{\TYP{EquivalenceRelation}\llbracket{}T,=\rrbracket}; with the
original declaration such a type is also an
\EXP{\TYP{EquivalenceRelation}\llbracket\TYP{Object},\sequiv\rrbracket},
and the two differ in a type argument. \TYP{Tuple}'s declaration is not
refused by the rule, since no tuple type extends a trait, but its text
follows \TYP{Object}'s. The replacement is the library's declaration.
\item[Effect] The properties of \EXP{\sequiv} are those stated in trait
\TYP{Any}; they are no longer restated at \TYP{Object} through the
algebra chapter.
\item[Original text] \nolinkurl{Specification-1.0-frozen/basic-lib/objects.tex}
(whose lines are two fewer than this document's source), line 82:
\begin{quote}
In \TYP{Object} this property is asserted abstractly using the algebraic constraints of \chapref{lib:algebraic-constraints}.
\end{quote}
lines 125--139, the declaration kept above, marked ``Not allowed'', after:
\begin{quote}
The \TYP{Object} type does not
have any additional methods, but uses algebraic constraints (see
\chapref{lib:algebraic-constraints}) to describe the properties of the
existing operators more abstractly.
\end{quote}
and lines 145--156:
\begin{quote}
As with the type \TYP{Object}, the type \TYP{Tuple} uses algebraic constraints to describe its existing methods.
\begin{Fortress}
\(\KWD{trait}\mskip 4mu plus 4mu\TYP{Tuple} \KWD{extends} \{\,\TYP{Any}, \TYP{EquivalenceRelation}\llbracket\TYP{Tuple},\sequiv\rrbracket, \TYP{IdentityOperator}\llbracket\TYP{Tuple}\rrbracket\,\}\)\\
{\tt~~~~}\pushtabs\=\+\(    \KWD{excludes} \{\,\TYP{Object}\,\}\)\-\\\poptabs
{\tt~~}\pushtabs\=\+\(  \VAR{toString}()\COLON \TYP{String}\)\-\\\poptabs
\(\KWD{end}\)
\end{Fortress}
\end{quote}
\item[Route C] By the revival's reading, not measured, the original
declaration of \TYP{Object} is one of the two passages that would read
differently under route~C, the numeric chapters being the other.
\end{description}

\subsection{The internal document}
\seclabel{revival-internal}

\begin{description}
\item[Affected sections] In the appendix ``Internal Document'', which
follows this one in a draft build: ``Disjunctive Constraints'',
``Weirdness with Multiple Inheritance and Polymorphism'' and
``Conditional Subtyping''.
\item[Change] One box after each: the first program's premise is not
allowed; the second is the case the rule answers; the third declaration is
not allowed for the reason the covariance example is not.
\item[Rationale] Instantiation exclusion.
\item[Effect] None; the internal document is printed only in a draft
build.
\item[Original text] Unchanged.
\item[Route C] The same under route~C.
\end{description}

\subsection{Passages not yet revised}
\seclabel{revival-notrevised}

The number chapters, \chapref{lib:numbers-basic} (whose integers are \secref{basic-integers})
and \chapref{lib:numbers-advanced}, are not yet revised.
Their nested tower, in which each integer type is a subtype of the rational
one above it and each carries its algebra at its own type, and the
conditional self-extension of \TYP{RationalQuantity}, are refused by
instantiation exclusion, as the compiled type checker refuses the
library's nested tower today.
They will be revised together with the reorganization of the library's
numeric types that route~A decides, and not before, because a
specification ahead of its library is as much an open discrepancy as one
behind it.
Until then the rule refuses the tower these chapters describe.

The calculi of \appref{calculi} predate the rule and are not yet revised:
the calculus with \KWD{where} clauses (\appref{where-calculus}) lets a
\KWD{where}-clause variable occur in a supertype, the calculus with an
acyclic type hierarchy admits a trait that extends an instantiation of
itself at its parameter's bound, and the basic calculus states that all of
its valid programs are valid Fortress programs, while it admits an object
that extends two instantiations of a parameterized trait declaring no
method.
How they are to be revised awaits a decision of the revival.

\subsection{Route C, the alternative not taken}
\seclabel{revival-routec}

The team's 2012 patent applications (US~8,843,887 and US~8,898,632;
Chase, Steele, Naden, Hilburn, and Luchangco) relax the rule for a
self-typed parameterized trait: a type may be a subtype of several of its
instantiations when they form a chain.
The revival built that rule whole as an experiment and measured it.
It closes the unsoundness that instantiation exclusion guards against only
together with a rule on return types over every level of a chain and a
change in how a call chooses among generic methods specialised to
different instantiations: 188 lines of code in five files of the compiled
type checker and code generator.
Reversing route~A for route~C would also need a method per level of a
chain in the code generator, work in the interpreter, a way past a stack
overflow of the type checker on the self-typing clauses, and seven library
declarations, besides undoing the reorganization of the library's numeric
types.
Route~C rescues none of the \KWD{where}-clause examples of
\secref{where-clauses} nor the \TYP{Nothing} of \secref{convenience-types},
each of which places one type below infinitely many instantiations, which
no chain contains.
Under route~C the numeric chapters could keep their nested tower, and the
original text of those chapters, kept unrevised in this document for now
and in the Working Draft of February 2011, is route~C's.
The record of the experiment is
\nolinkurl{explorations/reviews/mie-probes/route-c-experiment.md} in the
revival's repository.
\fussy
